\section{INTRODUÇÃO}

As buscas por animais perdidos dependem da divulgação do caso e da colaboração de pessoas que possam fornecer informações sobre o animal. Esse problema também foi observado em um estudo realizado nos Estados Unidos, no qual cerca de 15\% dos tutores perderam um cão ou um gato em um período de cinco anos (Weiss; Slater; Lord, 2012). Quando os avisos e relatos circulam por diferentes canais, acompanhar as atualizações pode se tornar difícil.

Para compreender melhor essas dificuldades, este trabalho realizou uma pesquisa exploratória com participantes da comunidade e analisou sistemas similares. As pessoas consultadas apontaram, principalmente, a perda de visibilidade das publicações e a dispersão das informações entre vários grupos. Também demonstraram incômodo com a exposição de contatos pessoais e interesse pela visualização dos casos em mapa. Esses resultados contribuíram para identificar necessidades de divulgação, organização das informações sobre os animais e registro de avistamentos.

Com base nessas necessidades, este trabalho apresenta o WeFIND, uma aplicação móvel colaborativa que apoia a divulgação e a organização de buscas por animais perdidos e encontrados.

O objetivo geral é desenvolver uma aplicação móvel colaborativa que apoie tutores e a comunidade na divulgação e na busca por animais perdidos e encontrados. Os objetivos específicos consistem em identificar as dificuldades relatadas por tutores e protetores na divulgação e na busca, analisar sistemas similares, definir os requisitos e modelar o sistema em UML (\textit{Unified Modeling Language}), projetar a interface, implementar o aplicativo com React Native e Supabase, testar o mecanismo de correspondência entre anúncios e avaliar a usabilidade do aplicativo por meio de tarefas e da Escala de Usabilidade do Sistema (SUS), aplicada em tradução para o português.

A relevância social do trabalho está em estudar e avaliar uma ferramenta para organizar a colaboração comunitária durante as buscas. A proposta reúne conhecimentos de Engenharia de Software, programação para dispositivos móveis, bancos de dados e design de interface, desenvolvidos no Curso Superior de Tecnologia em Sistemas para Internet do IFSul.

Metodologicamente, trata-se de uma pesquisa aplicada, de caráter exploratório-descritivo, com abordagens quantitativa e qualitativa. Os procedimentos são apresentados no Capítulo 3.

O trabalho está estruturado em nove capítulos, sendo esta introdução o primeiro deles. O Capítulo 2 apresenta o contexto dos animais perdidos, as formas de busca e o papel da colaboração e da localização geográfica. O Capítulo 3 descreve a metodologia, os resultados da pesquisa de campo, a análise de sistemas similares e o processo de desenvolvimento. O Capítulo 4 apresenta os requisitos e a modelagem do aplicativo. O Capítulo 5 descreve o WeFIND, seus fluxos e suas regras de negócio. O Capítulo 6 explica a arquitetura, as tecnologias, o banco de dados e o algoritmo de correspondência. O Capítulo 7 apresenta as decisões de interface e as principais telas. O Capítulo 8 descreve as verificações funcionais, a análise estática de segurança, o teste do mecanismo de correspondência e a avaliação de usabilidade, com seus resultados e limitações. Por fim, o Capítulo 9 reúne as considerações finais, as dificuldades encontradas e as possibilidades de trabalhos futuros.
